\documentclass[aspectratio=169]{beamer}

% Packages
\usepackage[utf8]{inputenc}
\usepackage{graphicx}
\usepackage{booktabs}
\usepackage{hyperref}
\usepackage{listings}
\usepackage{xcolor}

% Theme
\usetheme{Madrid}
\usecolortheme{default}
\useinnertheme{circles}

% Custom Colors (Chapar theme concept)
\definecolor{ChaparDarkBlue}{RGB}{20, 50, 90}
\definecolor{ChaparLightBlue}{RGB}{60, 120, 180}
\definecolor{ChaparAccent}{RGB}{220, 80, 40}

\setbeamercolor{palette primary}{bg=ChaparDarkBlue,fg=white}
\setbeamercolor{palette secondary}{bg=ChaparLightBlue,fg=white}
\setbeamercolor{palette tertiary}{bg=ChaparDarkBlue,fg=white}
\setbeamercolor{palette quaternary}{bg=ChaparDarkBlue,fg=white}
\setbeamercolor{structure}{fg=ChaparDarkBlue} % itemize, enumerate, etc
\setbeamercolor{section in toc}{fg=ChaparDarkBlue} % TOC sections
\setbeamercolor{alerted text}{fg=ChaparAccent}

% Title Page Information
\title[Chapar Project]{Chapar: Project Presentation}
\subtitle{Workshop Template}
\author[Author Name]{Your Name}
\institute[Institution]{Your Institution or Organization \\ \medskip \textit{your.email@example.com}}
\date{\today}
\titlegraphic{\vspace{-1cm}} % Can add a logo here: \includegraphics[width=2cm]{logo.png}

% Code Listing Style
\lstdefinestyle{mystyle}{
    backgroundcolor=\color{black!5},   
    commentstyle=\color{ChaparAccent},
    keywordstyle=\color{ChaparDarkBlue}\bfseries,
    numberstyle=\tiny\color{black!50},
    stringstyle=\color{ChaparLightBlue},
    basicstyle=\ttfamily\footnotesize,
    breakatwhitespace=false,         
    breaklines=true,                 
    captionpos=b,                    
    keepspaces=true,                 
    numbers=left,                    
    numbersep=5pt,                  
    showspaces=false,                
    showstringspaces=false,
    showtabs=false,                  
    tabsize=2
}
\lstset{style=mystyle}

\begin{document}

% Title Slide
\begin{frame}
    \titlepage
\end{frame}

% Table of Contents
\begin{frame}{Outline}
    \tableofcontents
\end{frame}

%---------------------------------------------------------
\section{Introduction}
%---------------------------------------------------------

\begin{frame}{Introduction to Chapar}
    \begin{itemize}
        \item Provide a brief overview of the Chapar project.
        \item What problem does it solve?
        \item Why is it relevant now?
    \end{itemize}
    \vspace{0.5cm}
    \begin{block}{Mission}
        To provide an efficient and reliable solution for [insert key capability].
    \end{block}
\end{frame}

%---------------------------------------------------------
\section{Architecture \& Design}
%---------------------------------------------------------

\begin{frame}{System Architecture}
    \begin{columns}
        \column{0.5\textwidth}
        \textbf{Key Components:}
        \begin{itemize}
            \item Component A: Description
            \item Component B: Description
            \item Component C: Description
        \end{itemize}
        
        \column{0.5\textwidth}
        % Uncomment and replace with your diagram
        % \includegraphics[width=\textwidth]{architecture.png}
        \begin{center}
            \textit{[Insert Architecture Diagram Here]}
        \end{center}
    \end{columns}
\end{frame}

%---------------------------------------------------------
\section{Key Features}
%---------------------------------------------------------

\begin{frame}[fragile]{Configuration Example} % [fragile] is needed when using code blocks
    \textbf{Easy to configure:}
    \vspace{0.3cm}
    \begin{lstlisting}[language=bash, caption={Example configuration snippet}]
# chapar configuration
repos:
  - name: my_repo
    url: https://example.com/repo
    branch: main
packages:
  - name: tool_a
    version: 1.2.0
    \end{lstlisting}
\end{frame}

%---------------------------------------------------------
\section{Usage \& Demo}
%---------------------------------------------------------

\begin{frame}{How to Use Chapar}
    \begin{enumerate}
        \item \textbf{Setup:} Initialize the environment.
        \item \textbf{Configure:} Define your parameters in \texttt{config.yaml}.
        \item \textbf{Execute:} Run the core pipeline.
    \end{enumerate}
    
    \vspace{0.5cm}
    \begin{alertblock}{Important Note}
        Always ensure your dependencies are correctly resolved before execution.
    \end{alertblock}
\end{frame}

%---------------------------------------------------------
\section{Future Work}
%---------------------------------------------------------

\begin{frame}{Roadmap}
    \begin{itemize}
        \item \textbf{Short-term:} Improve performance and add more documentation.
        \item \textbf{Medium-term:} Integrate with cloud providers.
        \item \textbf{Long-term:} Establish an open-source community around the tool.
    \end{itemize}
\end{frame}

%---------------------------------------------------------
\section{Conclusion}
%---------------------------------------------------------

\begin{frame}{Summary}
    \begin{itemize}
        \item Chapar simplifies complex deployments.
        \item Highly customizable via YAML configurations.
        \item Built with scalability in mind.
    \end{itemize}
    \vspace{1cm}
    \begin{center}
        \Large \textbf{Questions?}
    \end{center}
\end{frame}

\end{document}
